% Standalone problem document, generated from the adjacent README.md.
% Compile with XeLaTeX. Mathematical content is maintained in Markdown.
\documentclass[11pt,a4paper]{article}
\usepackage[fontsize=10.5pt]{scrextend}
\usepackage[margin=22mm,headheight=15pt,headsep=9mm,footskip=11mm]{geometry}
\usepackage{fontspec}
\setmainfont{texgyrepagella}[Extension=.otf,UprightFont=*-regular,BoldFont=*-bold,ItalicFont=*-italic,BoldItalicFont=*-bolditalic]
\setsansfont{texgyreheros}[Extension=.otf,UprightFont=*-regular,BoldFont=*-bold,ItalicFont=*-italic,BoldItalicFont=*-bolditalic]
\setmonofont{texgyrecursor}[Extension=.otf,UprightFont=*-regular,BoldFont=*-bold,ItalicFont=*-italic,BoldItalicFont=*-bolditalic,Scale=MatchLowercase]
\usepackage{amsmath,amssymb}
\usepackage{unicode-math}
\setmathfont{texgyrepagella-math.otf}
\usepackage{xcolor,microtype,enumitem,needspace,fancyhdr}
\usepackage{bookmark}
\definecolor{ink}{HTML}{17344B}
\definecolor{accent}{HTML}{197D80}
\definecolor{muted}{HTML}{596773}
\hypersetup{colorlinks=true,linkcolor=accent,urlcolor=accent,pdftitle={AV-03 - Polynomial-time
solution under the regularity
promise},pdfauthor={Open Problems in Numerical Linear Algebra}}
\urlstyle{same}
\setlength{\parindent}{0pt}
\setlength{\parskip}{5pt plus 1pt minus 1pt}
\linespread{1.04}
\setlength{\emergencystretch}{3em}
\widowpenalty=10000
\clubpenalty=10000
\displaywidowpenalty=10000
\allowdisplaybreaks[1]
\setlist{nosep,leftmargin=1.5em}
\providecommand{\tightlist}{\setlength{\itemsep}{0pt}\setlength{\parskip}{0pt}}
\providecommand{\pandocbounded}[1]{#1}
\pagestyle{fancy}
\fancyhf{}
\fancyhead[L]{\sffamily\footnotesize\color{muted}OPEN PROBLEMS IN NUMERICAL LINEAR ALGEBRA}
\fancyhead[R]{\sffamily\bfseries\footnotesize\color{accent}AV-03}
\fancyfoot[L]{\parbox[b]{0.9\textwidth}{\sffamily\fontsize{8}{10}\selectfont\color{muted}Editorial ratings; a bounded search cannot certify that a problem remains unsolved.\\Literature check: 2026-09-10}}
\fancyfoot[R]{\sffamily\footnotesize\color{muted}\thepage}
\renewcommand{\headrulewidth}{0.3pt}
\renewcommand{\headrule}{\hbox to\headwidth{\color{accent}\leaders\hrule height \headrulewidth\hfill}}
\makeatletter
\renewcommand{\section}{\Needspace{5\baselineskip}\@startsection{section}{1}{0pt}{12pt plus 2pt minus 2pt}{4pt}{\sffamily\bfseries\large\color{ink}}}
\renewcommand{\subsection}{\Needspace{4\baselineskip}\@startsection{subsection}{2}{0pt}{9pt plus 1pt minus 1pt}{3pt}{\sffamily\bfseries\normalsize\color{ink}}}
\makeatother
\setcounter{secnumdepth}{0}
\begin{document}
{\sffamily\small\color{muted}Intervals and absolute value equations\par}
\vspace{3pt}
{\raggedright\hyphenpenalty=10000\exhyphenpenalty=10000\sffamily\bfseries\fontsize{21}{24}\selectfont\color{ink}Polynomial-time
solution under the regularity promise\par}
\vspace{7pt}
{\sffamily\small\textbf{Difficulty:} extreme\quad\textbf{Importance:} broadly
interesting\par}
{\sffamily\small\bfseries\color{accent}Status: Open\par}
\vspace{4pt}
{\color{accent}\hrule height 0.8pt}
\vspace{6pt}
\textbf{Area:} algorithms for piecewise linear systems and
complementarity

\textbf{Rating rationale:} Extreme reflects the longstanding
polynomial-time barrier for P-matrix complementarity in an equivalent
regular AVE form; broad impact spans optimization, complexity, and
piecewise linear systems.

\subsection{Context and notation}\label{context-and-notation}

Absolute values and vector inequalities are componentwise. All
algorithmic questions use rational input in binary and the deterministic
Turing model. Input length includes the matrix dimensions and the bit
lengths of all rational coefficients.

For a real square matrix \(A\), define the diagonal perturbation family

\[
\mathcal D(A)=\{A-\operatorname{diag}(d):d\in[-1,1]^n\}.
\]

Call this family \textbf{regular} when every member is nonsingular. Only
diagonal entries vary. This is precisely the entrywise interval matrix
\([A-I_n,A+I_n]\).

\subsection{Problem statement}\label{problem-statement}

Is there a deterministic algorithm and a polynomial \(p\) such that, for
every \(n\ge1\), rational \(A\in\mathbb Q^{n\times n}\) and
\(b\in\mathbb Q^n\) with regular \(\mathcal D(A)\), the algorithm
outputs the exact rational vector \(x\) satisfying

\[
Ax-|x|=b
\]

within \(p(\ell)\) bit operations, where \(\ell\) is the binary input
length?

The promise guarantees existence and uniqueness. No certificate of
regularity is supplied or required, and behavior on inputs violating the
promise is unrestricted. Polynomial dependence on coefficient bit length
is allowed; this does not ask for a strongly polynomial algorithm.

\newpage
\subsection{References}\label{references}

Milan Hladík, Hossein Moosaei, Fakhrodin Hashemi, Saeed Ketabchi, and
Panos M. Pardalos,
\href{https://doi.org/10.1007/s10589-025-00717-5}{\emph{An overview of
absolute value equations: from theory to solution methods and
challenges}}, Computational Optimization and Applications \textbf{93}
(2026), 435--488, §2.4.2, paragraph following conditions (15)--(16);
§2.2 gives the complementarity connection.

This is the AVE formulation of the polynomial-time P-matrix linear
complementarity problem; it is counted once. Michaela Borzechowski, John
Fearnley, Spencer Gordon, Rahul Savani, Patrick Schnider, and Simon
Weber, \href{https://doi.org/10.4230/LIPIcs.ICALP.2024.32}{\emph{Two
Choices Are Enough for P-LCPs, USOs, and Colorful Tangents}}, ICALP
2024, §1, explicitly retain that algorithmic question.

\subsection{Earlier status evidence ---
2026-09-08}\label{earlier-status-evidence-2026-09-08}

The survey appeared online August 8, 2025. Borzechowski et al.'s latest
\href{https://arxiv.org/abs/2402.07683}{arXiv version} is v2, May 21,
2024. Searches for ``P-matrix linear complementarity 2026
polynomial-time'' and ``P-LCP solved polynomial algorithm'' found no
solution. E.-Nagy and Végh's
\href{https://arxiv.org/html/2605.10701v2}{June 30, 2026 revision},
abstract and §6, gives an algorithm whose complexity also depends on an
optimized handicap number; it does not give the required polynomial
bound in input length alone.

\subsection{Audit update --- 2026-09-10}\label{audit-update-2026-09-10}

Rechecked the
\href{https://link.springer.com/article/10.1007/s10589-025-00717-5}{2025
survey}, §2.4.2, and
\href{https://arxiv.org/html/2605.10701v2}{E.-Nagy--Végh's 2026 handicap
reduction}. The latter's complexity also depends on its handicap
parameter, so it does not establish a bound polynomial solely in the
displayed input length. P-LCP/regular-AVE polynomial-time searches found
no such algorithm or matching resolution.
\end{document}
